\section{Lean definitions, checked results, and completion plan}
\label{sec:lean}

\subsection{Two different formalization levels}

The sources use Lean 4.32.0 and its bundled \texttt{Std} library.
They contain native definitions and kernel-checked arithmetic theorems,
not yet proof objects in ProveIt's first-order theory. Lean's metatheory
is stronger than iPRA even when the mathematical argument has a short
iPRA proof.

The arithmetic and search modules are \texttt{Arithmetic.lean} and
\texttt{FiniteSearch.lean}. The other modules are
\texttt{QuarticExample.lean} and \texttt{Specifications.lean}.
All four compile.
The audit module prints axiom dependencies and checks target types.
The general targets are definitions of propositions, not asserted
theorems or added axioms.

\subsection{Checked native mathematics}

\texttt{Arithmetic.lean} proves the nonzero-product bound, the
difference-of-squares identity, and the decisive gap statement
\[
 z^2=a^2+r,\quad r\ne0\quad\Longrightarrow\quad |a|\le|r|.
\]
It proves the quartic factorization and both sharp quartic bounds,
uniformly in $a,b,x,y$, and proves that no natural-number square lies
strictly between consecutive natural-number squares.

\texttt{FiniteSearch.lean} defines
\begin{lstlisting}
def scan (p : Nat -> Bool) (B : Nat) : List Nat :=
  (List.range (B+1)).filter p
\end{lstlisting}
and proves its exact membership characterization. It proves that a
complete finite list and a numerical bound are equivalent for predicates
on natural-number codes. It also proves the intentionally weaker
fact that every finite prefix is bounded.

\texttt{QuarticExample.lean} proves for all integers $x,y$ that
\[
\begin{split}
y^2=x^4+7x^2+1\quad\Longleftrightarrow\quad
&[x=0\ \wedge\ (y=1\vee y=-1)]\\
{}\vee{}&[(x=1\vee x=-1)\ \wedge\ (y=3\vee y=-3)].
\end{split}
\]
Its global coordinate bounds are proved before the finite case analysis.
This is not a bounded experiment promoted to a universal theorem.

The axiom audit contains only the standard Lean principles
\texttt{propext} and \texttt{Quot.sound}; the finite-prefix witness
requires none. No audited result depends on \texttt{Classical.choice},
\texttt{sorryAx}, or an added theorem axiom. This fact does not replace
the missing embedded iPRA derivations.

\subsection{Exact specifications and unproved targets}

The specification file defines polynomials as \texttt{List Int},
Horner evaluation, coefficient lookup with default zero, a coefficient
norm, and convolution coefficients. A certificate contains the finite
data $f,A,R,D$. The Boolean \texttt{certificateCheck} verifies positivity,
degree, nonzero conditions, and coefficient equality over a finite
support range. It contains no assumed height bound.

The bound functions are implemented directly:
\[
\begin{split}
\texttt{xBound}(c)&=\|A_{<\deg A}\|_1+\|R\|_1,\\
\texttt{yBound}(c)&=\sum_i |f_i|\max(1,\texttt{xBound}(c))^i.
\end{split}
\]
The file also defines the integer interval, rectangular filter, and
$X^{2m}+X+1$. The propositions \texttt{RungeBoundTarget},
\texttt{RungeEnumerationTarget}, and \texttt{GenusFamilyTarget}
are unproved in Lean. The corresponding mathematics is proved in
the paper. A concrete sample certificate is checked, and its finite
box evaluates to the expected four points.

For general affine systems the file defines monomials, exponent lists,
polynomial evaluation, a well-formedness test, an integral-point test,
and coordinate height. \texttt{PositiveSiegelTarget} and
\texttt{CountSiegelTarget} have different conclusions. Each accepts an
\texttt{admissible} predicate as a parameter. Implementing that
predicate so that it really certifies smooth positive-genus curves,
and proving its coverage, remain separate obligations. A parameterized
definition asserts no geometric theorem.

\subsection{The actual ProveIt interfaces}

At the pinned revision the namespace is
\texttt{LeanProofs.ConciseIPRA}. The inspected interfaces are
\begin{center}
\begin{tabular}{@{}p{0.35\linewidth}p{0.56\linewidth}@{}}
\toprule
Interface & Role\\
\midrule
\texttt{PR k} & Primitive recursive construction of arity $k$\\
\texttt{PTm n} & Arithmetic terms with $n$ free variables\\
\texttt{Fm PTm n} & First-order formulas\\
\texttt{Fm.QF} & Quantifier-free restriction\\
\texttt{PRAAx} & Arithmetic axioms and restricted induction\\
\texttt{Prf PRAAx Gamma phi} & Intuitionistic natural-deduction derivability\\
\texttt{PRA.qfInd} & Quantifier-free induction rule\\
\texttt{pra\_iff\_native} & Transfer to the concise native calculus\\
\bottomrule
\end{tabular}
\end{center}

The intended elementary result must ultimately have the type
\begin{lstlisting}
Prf PRAAx [] rungeFormula
\end{lstlisting}
where \texttt{rungeFormula} is the closed universal formula built
from the compiled \texttt{PR} functions and actual \texttt{Fm}
constructors. A host proposition with similar mathematical words
is insufficient. Once the derivation exists,
\texttt{pra\_iff\_native} transports it; the equivalence does not
create the missing derivation.

\subsection{Dependency-ordered completion tasks}

First prove that the finite polynomial checker implies the evaluated
identity. Establish the coefficient and integer-root bounds for
normalized lists. Combine them with \texttt{square\_gap} to prove
\texttt{RungeBoundTarget}. Prove membership in the integer interval
and nested filter to obtain \texttt{RungeEnumerationTarget}.
Instantiate the certificate uniformly in $m$ and apply the consecutive
square lemma to prove \texttt{GenusFamilyTarget}. Every mathematical
ingredient appears in Section~\ref{sec:runge-certificate}.

Next construct fixed primitive recursive codes for signed arithmetic,
arrays, sums, evaluation, the certificate test, both bounds, and
enumeration. Derive their equations inside \texttt{PRAAx}. Compile
the invariants of Section~\ref{sec:arithmetic-audit} to characteristic
equations, supply their \texttt{QF} witnesses, and use
\texttt{PRA.qfInd}. Assemble the universal derivation by natural
deduction and equality substitution, then specialize it to concrete
certificates and transport it to the concise calculus.

The accompanying \texttt{OBJECT\_PLAN.txt} expands these steps into
individual invariants and proof obligations. Completion requires
the actual object derivation, an axiom audit, and a clean build.
A native Lean build alone does not meet that requirement.

There is a finite implementation path for the elementary results.
The general Siegel dependencies in Section~\ref{sec:general-route}
are additional mathematical research tasks. Completing routine Lean
infrastructure is not promised to resolve them.

